\chapter{Происхождение Hubbard--Stratonovich odd auxiliary parent}
\label{ch:v10-h15-r2-hubbard-stratonovich-parent}

Предыдущий аудит выделил минимальное восьмимерное поле $A_\Sigma$ как
единственный положительный алгебраически устранимый кандидат. Проверим,
является ли Hubbard--Stratonovich преобразование происхождением gap или
лишь точной перезаписью уже известного оператора.

Пусть $Q=6Y$ на восьми секторах $\Sigma$. Формальная квадратичная форма
равна
\begin{equation}
 V_{\rm HS}=\frac12\Sigma^T49I\Sigma+A_\Sigma^TQ\Sigma
 +\frac12A_\Sigma^TA_\Sigma.
\end{equation}
Завершение квадрата даёт точное тождество
\begin{equation}
 \boxed{V_{\rm HS}=\frac12\|A_\Sigma+Q\Sigma\|^2
 +\frac12\Sigma^T(49I-Q^2)\Sigma.}
\end{equation}
Stationary equation фиксирует $A_\Sigma^\star=-Q\Sigma$. Triangular shift
имеет determinant $1$, auxiliary metric равна $I_8$ и также имеет
determinant $1$. Поэтому нормированное Gaussian integration не создаёт
field-dependent determinant и точно возвращает Schur complement $G_Y$.

Но это тождество начинается с target-loaded mixed bilinear
$A_\Sigma^TQ\Sigma$. Унаследованный source-free parent имеет Hessian
\begin{equation}
 H_0=\begin{pmatrix}49I&0\\0&I\end{pmatrix},
 \qquad \rank H_0=16,
\end{equation}
и после устранения $A_\Sigma$ оставляет только $49I$. Требуемый parent
\begin{equation}
 H_{\rm HS}=\begin{pmatrix}49I&Q\\Q&I\end{pmatrix}
\end{equation}
имеет rank $14$ и nullity $2$. Следовательно, $H_0$ и $H_{\rm HS}$ не
могут быть связаны обратимой заменой полей: congruence сохраняет rank.

Разность двух форм есть новый оператор
\begin{equation}
 \Delta H=\begin{pmatrix}0&Q\\Q&0\end{pmatrix}.
\end{equation}
Он имеет rank $16$ и inertia $(8_-,0_0,8_+)$. То есть HS-поле корректно
линеаризует уже присутствующий отрицательный $Q^2$-вклад, но не выводит
отсутствующий mixed source из universal mass и Gaussian measure.

\begin{theorem}[HS linearization без parent-origin]
Hubbard--Stratonovich преобразование с источником $Q\Sigma$ точно
эквивалентно effective gap $49I-Q^2$ и сохраняет нормированную меру.
Однако source-free inherited parent имеет другой rank, а добавляемый
cross-оператор имеет полный rank $16$. Поэтому преобразование является
переписыванием target-term, а не его физическим происхождением.
\end{theorem}

\begin{nogocriterion}[Gaussian identity не создаёт mixed source]
Нельзя считать $-Q^2$ выведенным из Hubbard--Stratonovich процедуры, если
bilinear $A_\Sigma^TQ\Sigma$ введён специально для получения этого члена.
Сначала должен быть независимо получен типизированный cross-оператор $Q$ в
исходном parent action.
\end{nogocriterion}

\begin{protocol}\begin{enumerate}
\item stationary auxiliary map rank: \textbf{$8$};
\item triangular shift determinant: \textbf{$1$};
\item Gaussian metric determinant: \textbf{$1$};
\item HS parent rank/nullity: \textbf{$14/2$};
\item inherited parent rank/nullity: \textbf{$16/0$};
\item inherited/required cross rank: \textbf{$0/8$};
\item new operator increment rank: \textbf{$16$};
\item increment inertia: \textbf{$(8_-,0_0,8_+)$};
\item formal/physical origin: \textbf{$4/4$ и $3/4$};
\item следующий гейт: \textbf{аудит кандидатов odd auxiliary cross-bilinear}.
\end{enumerate}\end{protocol}

Машинный сертификат сохранён в
\path{s2t_v10_particle_wrinkle_dislocation_callias_equal_charge_twist_m4_h15_r2_hubbard_stratonovich_odd_auxiliary_parent_origin_gate_results.json}.